\section{Problem Difficulty Analysis and Proposed Method}
\label{sec:method}

\subsection{Why this is hard}
\label{sec:difficulty}

Three things make in-place permutation of an N-D tensor harder than in-place
matrix transposition.

\subsubsection{The permutation is not a single transpose}

Cycle following works on a matrix because the index map has a closed form.
Position $p$ of an $m \times n$ matrix moves to
$(p \bmod n)\,m + \lfloor p/n \rfloor$, and the orbits of that map partition
the buffer. An arbitrary $\pi$ over $r$ axes induces an index map with no such
closed form. It can still be decomposed into disjoint
cycles, because every permutation can~\cite{dummit2003abstract}. But computing
those cycles means materialising the map. That is $\Theta(N)$ work and
$\Theta(N)$ state, which defeats the purpose.

\subsubsection{The decomposition is not unique, and the choices are not equal}

Any $\pi$ can be reached by a sequence of simpler moves, and there are many
such sequences. They are not interchangeable. Each move rewrites the whole
buffer, so a decomposition of $k$ steps pays roughly $k$ passes over memory.
A decomposition that is short in steps may still move more bytes than a longer
one, because what a step actually moves depends on how many positions it
leaves fixed. Choosing well is a search problem rather than a bookkeeping one.

\subsubsection{Memory access, not arithmetic, is the cost}

A permutation does no arithmetic. Every step is bound by how the memory system
responds to the access pattern that step produces. That response is strongly
non-linear in the stride. A step whose innermost contiguous run is long enough
to fill a cache line runs at a large multiple of the throughput of one whose
run is short. A cost model for this cannot be derived from operation counts.
It has to be calibrated against the hardware.

The search space is $r!$ orderings of the axes, with an edge between two
orderings for every move that connects them. An exhaustive comparison is out
of reach at the ranks that matter.

\subsection{Proposed method: adjacent box swaps}
\label{sec:boxes}

We propose decomposing $\pi$ into a sequence of \emph{adjacent box swaps}.

A \emph{box} is one or more neighbouring axes treated as a single unit. Those
axes are adjacent in a row-major layout, so their extents multiply into one
logical extent and the group behaves exactly like a single axis of that size.
A \emph{swap} exchanges two adjacent boxes.

The point of the construction is what a swap reduces to. Fix a split of the
axis list into a prefix, a left box, a right box and a suffix, with extents
$D_{pre}$, $\mathit{rows}$, $\mathit{cols}$ and $D_{post}$. Exchanging the two
boxes is exactly an in-place transpose of a logical
$\mathit{rows} \times \mathit{cols}$ matrix. It is performed independently for
each of the $D_{pre}$ leading positions, and each element of it is a
contiguous run of $D_{post}$ values. Every step of the decomposition is
therefore a problem we already know how to solve in place, and the
two-dimensional cycle-following result applies unchanged at every step.

$D_{post}$ also determines whether a step is fast. A large $D_{post}$ means
each moved element is a long contiguous run, which is what the memory system
rewards. When it is divisible by the vector width, the kernel moves four
values at a time instead of one.

\subsection{Choosing the decomposition by shortest path}
\label{sec:planning}

A permutation can always be reached by bubbling each axis into position with
single-axis swaps. Treating a run of adjacent axes as one box instead turns
the decomposition into a shortest-path problem over axis orderings. The search
is then free to spend one swap where the baseline spends several.

We build a graph whose vertices are orderings of the axes and whose edges are
legal box swaps. Each edge is weighted by the bytes that swap moves, divided
by a calibrated throughput. Dijkstra from the identity ordering to $\pi$ then
returns the cheapest decomposition under the cost model rather than merely a
correct one.

To measure what that freedom is worth we sample shape and permutation pairs
from the space defined in Section~\ref{sec:benchmark} and plan each of them
both ways. The sweep covers \planRanks{} ranks, \planRankLo{} through
\planRankHi{}, at \planElements{} elements throughout. It draws
\planShapesPerRank{} shapes per rank and up to \planPermsPerShape{}
permutations per shape, giving \planCases{} pairs and \planPlans{} plans. The
comparison is paired. Both planners see the same shape and the same
permutation, and every rank holds the element count fixed, so a difference in
bytes moved is a difference in the decomposition rather than in the tensor.

Across the sweep the search issues \planStepFactor{}x fewer swaps than the
adjacent baseline and moves \planByteFactor{}x fewer bytes at the median. It
reaches \planByteFactorMax{}x on the permutation where it helps most. The
margin widens with rank, as Figure~\ref{fig:steps} shows. At rank
\planRankHi{} the median permutation costs the baseline
\planStepsBaselineTop{} swaps and the search \planStepsSearchTop{}.
Table~\ref{tab:planners} gives the breakdown. \planBlockShare{}\% of the swaps
the search selects bundle more than one axis into a box. That is the move the
baseline cannot make, and it is where the difference comes from.

The auxiliary cost of a plan is its cycle metadata, and it does not scale with
the tensor. The largest plan in the sweep needs \planAuxWorst{} bytes, which
is \planAuxPercent{}\% of the data it permutes. The data itself is never
copied to a second buffer. That is the objective the whole construction exists
to serve.

% Table (tab:planners). Body: _generated/planners-table.tex, written by
% pipeline/tables/gen_table_planners.py from raw_data/plan_sweep/.
%
% The cells are generated. This wrapper -- the float, the caption, the label --
% is hand-written, and quotes the table's own macros so the caption cannot
% disagree with the rows underneath it.
%
% table*, not table: seven columns of data have a natural width of about 283pt
% and a single USENIX column is 241pt. tabular* silently overflows rather than
% warning, so this was overprinting the neighbouring column. The full measure
% fits it at \small with room to spare; shrinking the type instead would have
% meant scriptsize, which is 8pt at this base size.
\begin{table*}[!t]
\centering
\caption{\capheadline{What the box search saves over the adjacent baseline, by
rank.} \planTabCases{} shape and permutation pairs across \planTabRanks{}
ranks, sampled from the shape space of Equation~(\ref{eq:shapes}) and planned
by both strategies at a fixed tensor size. The bytes-moved advantage grows
with rank, from \planTabWorstX{}x where a permutation is a single transpose to
\planTabBestX{}x at rank~\planTabBestRank{}.}
\label{tab:planners}
% GENERATED by pipeline/tables/gen_table_planners.py -- do not edit by hand.
% Regenerate after any change to raw_data/plan_sweep/ (experiments/plan_sweep/run.py).
\small
\begin{tabular*}{\linewidth}{@{\extracolsep{\fill}}l r r r rr r r@{}}
\toprule
 & & & & \multicolumn{2}{c}{\textbf{Swaps}} & & \\
\cmidrule(lr){5-6}
\textbf{Rank} & \textbf{$\Omega(N,r)$} & \textbf{Shapes} & \textbf{Pairs} & \textbf{Adj.} & \textbf{Search} & \textbf{Bytes} & \textbf{Aux (B)} \\
\midrule
2 & 96 & 12 & 12 & 1 & 1 & 1.0x & 304 \\
3 & 1{,}944 & 12 & 48 & 2 & 1 & 1.5x & 592 \\
4 & 19{,}200 & 12 & 48 & 3 & 2 & 2.0x & 3{,}728 \\
5 & 123{,}750 & 12 & 48 & 5 & 2 & 2.1x & 2{,}296 \\
6 & 598{,}752 & 12 & 48 & 8 & 2.5 & 2.9x & 720 \\
\bottomrule
\end{tabular*}

\vspace{1.5pt}
\parbox{\linewidth}{\fontsize{6}{6.8}\selectfont\raggedright $\Omega(N,r)$ is the size of the shape space of Equation~(\ref{eq:shapes}) at this rank. \textbf{Shapes} are drawn from it uniformly and \textbf{Pairs} counts (shape, permutation) cases. Swaps and bytes are medians over the pairs in that row. \textbf{Bytes} is the median per-pair ratio of bytes moved by the adjacent baseline to bytes moved by the search, so higher is better. \textbf{Aux (B)} is the largest cycle-metadata footprint any one search plan in the row needed, in bytes.}
%
\end{table*}

% Figure (fig:steps). Body: _generated/steps-fig.tex, written by
% pipeline/figures/plot_planner_steps.py from raw_data/plan_sweep/.
%
% The marks are generated TikZ, drawn by this document's own build, so the
% labels are set in the paper's fonts at the sizes the renderer asked for.
% Nothing here scales the picture: a \resizebox around a chart is what makes
% one paper's figures come out with labels at four different sizes.
%
% The node-and-rectangle pair frames the chart at whatever size its bounding
% box declares. Colours defined inside the generated file are scoped to it,
% so the frame uses black rather than one of the chart's own.
\begin{figure}[!t]
\centering
\begin{tikzpicture}
\node[inner sep=2pt] (figsteps) {% GENERATED by pipeline/figures/plot_planner_steps.py -- do not edit by hand.
% Regenerate after any change to raw_data/plan_sweep/ (experiments/plan_sweep/run.py).
\begingroup%
\definecolor{pink}{HTML}{1C1B18}%
\definecolor{pmid}{HTML}{6B6A65}%
\definecolor{pgrid}{HTML}{C9C7C0}%
\definecolor{pblue}{HTML}{2A78D6}%
\definecolor{porange}{HTML}{EB6834}%
\begin{tikzpicture}[x=1pt,y=1pt]%
\useasboundingbox (0,0) rectangle (241.00,132.00);%
\draw[pgrid,line width=0.4pt] (30.00,26.00) -- (233.00,26.00);%
\node[anchor=east,text=pink,font=\fontsize{5.5}{6.5}\selectfont] at (28.00,26.00) {0};%
\draw[pgrid,line width=0.4pt] (30.00,44.80) -- (233.00,44.80);%
\node[anchor=east,text=pink,font=\fontsize{5.5}{6.5}\selectfont] at (28.00,44.80) {2};%
\draw[pgrid,line width=0.4pt] (30.00,63.60) -- (233.00,63.60);%
\node[anchor=east,text=pink,font=\fontsize{5.5}{6.5}\selectfont] at (28.00,63.60) {4};%
\draw[pgrid,line width=0.4pt] (30.00,82.40) -- (233.00,82.40);%
\node[anchor=east,text=pink,font=\fontsize{5.5}{6.5}\selectfont] at (28.00,82.40) {6};%
\draw[pgrid,line width=0.4pt] (30.00,101.20) -- (233.00,101.20);%
\node[anchor=east,text=pink,font=\fontsize{5.5}{6.5}\selectfont] at (28.00,101.20) {8};%
\draw[pgrid,line width=0.4pt] (30.00,120.00) -- (233.00,120.00);%
\node[anchor=east,text=pink,font=\fontsize{5.5}{6.5}\selectfont] at (28.00,120.00) {10};%
\node[anchor=north,text=pink,font=\fontsize{5.5}{6.5}\selectfont] at (46.92,23.50) {2};%
\node[anchor=north,text=pink,font=\fontsize{5.5}{6.5}\selectfont] at (89.21,23.50) {3};%
\node[anchor=north,text=pink,font=\fontsize{5.5}{6.5}\selectfont] at (131.50,23.50) {4};%
\node[anchor=north,text=pink,font=\fontsize{5.5}{6.5}\selectfont] at (173.79,23.50) {5};%
\node[anchor=north,text=pink,font=\fontsize{5.5}{6.5}\selectfont] at (216.08,23.50) {6};%
\draw[pgrid,line width=0.6pt] (30.00,26.00) -- (233.00,26.00);%
\draw[pgrid,line width=0.6pt] (30.00,26.00) -- (30.00,120.00);%
\draw[porange,line width=0.9pt] (46.92,35.40) -- (89.21,44.80) -- (131.50,54.20) -- (173.79,73.00) -- (216.08,101.20);%
\fill[white] (45.42,33.90) rectangle (48.42,36.90);%
\draw[porange,line width=0.7pt] (45.42,33.90) rectangle (48.42,36.90);%
\fill[white] (87.71,43.30) rectangle (90.71,46.30);%
\draw[porange,line width=0.7pt] (87.71,43.30) rectangle (90.71,46.30);%
\fill[white] (130.00,52.70) rectangle (133.00,55.70);%
\draw[porange,line width=0.7pt] (130.00,52.70) rectangle (133.00,55.70);%
\fill[white] (172.29,71.50) rectangle (175.29,74.50);%
\draw[porange,line width=0.7pt] (172.29,71.50) rectangle (175.29,74.50);%
\fill[white] (214.58,99.70) rectangle (217.58,102.70);%
\draw[porange,line width=0.7pt] (214.58,99.70) rectangle (217.58,102.70);%
\draw[pblue,line width=0.9pt] (46.92,35.40) -- (89.21,35.40) -- (131.50,44.80) -- (173.79,44.80) -- (216.08,49.50);%
\fill[white] (46.92,35.40) circle (1.5pt);%
\draw[pblue,line width=0.7pt] (46.92,35.40) circle (1.5pt);%
\fill[white] (89.21,35.40) circle (1.5pt);%
\draw[pblue,line width=0.7pt] (89.21,35.40) circle (1.5pt);%
\fill[white] (131.50,44.80) circle (1.5pt);%
\draw[pblue,line width=0.7pt] (131.50,44.80) circle (1.5pt);%
\fill[white] (173.79,44.80) circle (1.5pt);%
\draw[pblue,line width=0.7pt] (173.79,44.80) circle (1.5pt);%
\fill[white] (216.08,49.50) circle (1.5pt);%
\draw[pblue,line width=0.7pt] (216.08,49.50) circle (1.5pt);%
\node[anchor=north,text=pink,font=\fontsize{6}{7}\selectfont] at (131.50,15.00) {tensor rank};%
\node[anchor=south,rotate=90,text=pink,font=\fontsize{6}{7}\selectfont] at (13.00,73.00) {median swaps per permutation};%
\draw[porange,line width=0.9pt] (38.46,116.24) -- (47.46,116.24);%
\node[anchor=west,text=pink,font=\fontsize{5.5}{6.5}\selectfont] at (49.46,116.24) {adjacent baseline};%
\draw[pblue,line width=0.9pt] (38.46,108.74) -- (47.46,108.74);%
\node[anchor=west,text=pink,font=\fontsize{5.5}{6.5}\selectfont] at (49.46,108.74) {box search};%
\end{tikzpicture}%
\endgroup%
};
\draw[black,line width=0.5pt,rounded corners=2pt]
  (figsteps.south west) rectangle (figsteps.north east);
\end{tikzpicture}
\caption{\capheadline{Median swaps per permutation against tensor rank.}
\stepsFigRanks{} ranks, both planners on the same shapes. The baseline's cost
climbs with the number of adjacent transpositions a permutation needs. The
search stays flat. At rank~\stepsFigTopRank{} it ends at
\stepsFigTopSearch{} swaps against \stepsFigTopBaseline{}.}
\label{fig:steps}
\end{figure}


\subsection{Feasibility}
\label{sec:feasibility}

A working prototype exists, and the numbers above come from it rather than
from an estimate. The planner is implemented in \texttt{src/\_plan.py} and the
cycle-following kernel in \texttt{src/\_kernels/box\_kernel.cu}, with both a
scalar and a vectorised path. What remains is calibration and evaluation
rather than construction. That is why the timeline in
Section~\ref{sec:milestones} is weighted towards measurement.

Edge weight is bytes moved divided by a calibrated throughput $\mu$, currently
\costMu{}~GB/s. Two thresholds gate the search. $\tau$ is the smallest
$D_{post}$ at which an edge is worth considering. $D_{post}^{*}$ is the point
at which box following beats a cycle splitter. Both currently sit at
\costTau{} and \costDPostStar{}, so the pruning they describe is
\costPruning{}. Re-deriving them on the evaluation hardware is a milestone in
its own right.

\placeholder{One known gap to close before the final report. During the search
\texttt{estimated\_cycle\_info} reports no fixed points, so every edge weighs
the same and the search minimises swap count rather than bytes. Exact counts
are filled in afterwards, so the reported numbers are accurate for the plans
chosen, but they are not what the search optimised. Quantify the resulting
gap.}

\subsection{Why this improves on existing work}
\label{sec:improvement}

We compare against prior work along the dimensions that matter for this
problem. These are the number of full-tensor passes, the auxiliary memory a
method needs, how sensitive its cost is to the tensor shape, and whether it
supports arbitrary N-D permutations.

\textbf{Data movement.} Most GPU-oriented approaches trade extra data movement
for regular, well-behaved
kernels~\cite{catanzaro2014decomposition,gomezluna2016inplace,wu2025eithot,cheng2025ittpd}.
They may read and write the tensor several times even when the permutation
itself requires moving each element once. Cycle following is the exception,
since it moves each element exactly once~\cite{gustavson2012cycleleader}, but
the original formulation is built around individual cycles and does not give
an efficient GPU implementation for arbitrary N-D permutations. Our proposal
sits with the first group and accepts multiple passes. Its claim is that it
takes fewer of them than a fixed decomposition rule does, and the sweep above
is the evidence.

\textbf{Auxiliary memory.} Existing methods size their temporary storage from
the tensor shape or from the decomposition strategy. Catanzaro et
al.~\cite{catanzaro2014decomposition} need storage related to the larger
dimension. Gomez-Luna et al.~\cite{gomezluna2016inplace} may pad an axis to
obtain a usable factorization, so the workspace depends on the input.
\textsc{ITTPD}~\cite{cheng2025ittpd} exposes a memory-budget parameter, but
uses it to set how many batch instances its primitive kernels process. Our
auxiliary state is the cycle metadata of the selected steps, which the sweep
measures at \planAuxPercent{}\% of the tensor in the worst case.

\textbf{Shape sensitivity.} Several of these designs rely on cache locality to
reach peak
throughput~\cite{catanzaro2014decomposition,gomezluna2016inplace,cheng2025ittpd}.
A shape that yields short or irregular accesses costs them disproportionately.
A search over decompositions attacks this directly, because the cost model
weights an edge by the access pattern that edge produces. Whether that
translates into more stable measured runtime is an open question and is a
deliverable of the evaluation.

\textbf{Permutation coverage.} Catanzaro et
al.~\cite{catanzaro2014decomposition}, Gustavson et
al.~\cite{gustavson2012cycleleader} and Gomez-Luna et
al.~\cite{gomezluna2016inplace} define their kernels for two-dimensional
permutations or single axis swaps. Data layout transformation systems apply
the same machinery to strided layouts~\cite{sung2012pttwac}. None of them
addresses an arbitrary permutation of an arbitrary rank directly.
\textsc{EITHOT}~\cite{wu2025eithot} and
\textsc{ITTPD}~\cite{cheng2025ittpd} do, and \textsc{ITTPD} decomposes the
permutation before transposing it, which is the same overall shape of idea as
ours.

\textbf{What is different here.} Against \textsc{ITTPD} in particular, our
proposal differs in two ways, and these are the claims we intend to defend.

\begin{enumerate}
\item \textbf{Boxes, not axes.} A step may exchange two \emph{groups} of
  adjacent axes rather than two axes. This strictly enlarges the set of
  reachable decompositions. The measurements above show the search takes that
  option on \planBlockShare{}\% of the steps it selects, so it is not a
  generalisation that goes unused.
\item \textbf{Search under a calibrated cost model, not a fixed rule.} The
  decomposition is chosen by shortest path over a graph whose weights are
  measured throughput. The planner therefore adapts to the hardware it is
  calibrated on, instead of encoding one machine's behaviour as a heuristic.
\end{enumerate}

\placeholder{Relationship to Cycle Splitter, our group's other approach to the
same problem, needs stating here with a citation. Cycle Splitter moves each
element exactly once and bounds auxiliary memory explicitly. Box following
accepts multiple passes in exchange for regular, coalesced kernels, and wins
where $D_{post}$ is large enough. The constant $D_{post}^{*}$ in
\texttt{src/\_plan.py} is the routing gate between the two, so the comparison
is not rhetorical. Deriving it is a milestone.}

\placeholder{A direct head-to-head against \textsc{EITHOT} and \textsc{ITTPD}
on their reported shapes is the strongest evidence for these claims, and is
not yet run. It is the deliverable of the evaluation milestone.}
